\section{Implementation for ResNet-50}
\label{app:implementation-resnet50}

Both pruning variants use the shared objective in
Appendix~\ref{app:shared-loss}, with normalization over remaining channels
and sampled ST training masks.

\subsection{Iterative Pruning}
\label{sec:resnet-iterative-implementation}
Each trajectory uses the iterative composition in Section~\ref{sec:iterative-method}, with cycle
and phase budgets in Table~\ref{tab:pruning-schedules}.
Every search and recovery phase starts with newly initialized random
weights; search weights are not inherited. The trajectories use threshold or budget-only selection, as listed
in Table~\ref{tab:method-hyperparameters}. Only the cumulative removed-channel set is carried forward.
The parameter counts and test accuracies of all ten recovery models in
each series are listed in Table~\ref{tab:resnet-notebook-trajectories}.
% Initialization and phase budgets confirmed by the author, 2026-11-11.
% The source notebook does not establish the remaining historical gate/controller
% settings or a possible additional cap for the two threshold trajectories.

\subsection{One-Shot Pruning}
\label{sec:resnet-cifar-implementation}
The single-pass composition uses Algorithms~\ref{alg:search-phase}--\ref{alg:recovery-phase},
with history-feasible selection and budgets in Table~\ref{tab:pruning-schedules}.
Recovery inherits the surviving search weights, recalibrates BatchNorm
on 200 training batches and starts a fresh optimizer. We test the
best validation recovery checkpoint with frozen weights and BatchNorm.

We use a CIFAR-adapted ResNet-50~\cite{he2015residual}
(23.547M parameters), with a $3\times3$ stride-1 stem and no
max-pooling. The eight repeated runs share random initialization
(seed 42); the separately trained dense model supplies only the
validation reference. All parameter groups start at learning rate
$10^{-3}$; entropy regularization and label smoothing are disabled.

\paragraph{Gates and regularization.}
Each of the 16 bottlenecks has two independently gated internal
widths, after the first and second convolutions' BatchNorm and ReLU:
32 groups, with no output gates. Search uses sampled ST masks;
evaluation thresholds opening probabilities. The initial opening
probability and evaluation threshold are listed in
Table~\ref{tab:method-hyperparameters}. Loss components and normalization
are defined in Appendix~\ref{app:shared-loss}.

\paragraph{Sparsity control and selection.}
The controller averages paired, epoch-matched validation gaps.
The averaging window, tolerances, warmup, initial coefficient, clipping
bounds and automatic log-step are listed in
Table~\ref{tab:method-hyperparameters}. This configuration does not use
the iterative closure-rate boost.
From epochs 31--60 we select the smallest feasible physical model,
with a 0.001 retention floor, equivalent here to retaining at least
one channel per group; fallback maximizes validation accuracy.
Candidates are scored before additional BatchNorm calibration.

Earlier S1/S2 tolerance sweeps used retention floors of
0.001/0.08. S1 feasible candidates came from epochs 1--60, with
fallback from 31--60; both S2 windows were 31--60. The additional
mean represents two selected repeats of the configuration above.
% Sources and historical-series distinctions: EXPERIMENTS.md;
% saved resnet-implementation and algorithm-clarification research.
